\section{Introduction}\label{introduction}

Agentic LLMs now attempt long, multi-step tasks, but a human still babysits: watching, correcting,
telling the agent the next step. Remove the human and one failure dominates --- the agent declares
success it never checked. This is worse than an ordinary error: it is a \emph{silent, corrupting} one,
because the very signal that something is wrong (the human\textquotesingle s glance) has been removed. Unattended
autonomy is therefore gated not by capability but by \emph{trust}.

Existing fixes do not close this gap. Self-correction (Reflexion, Self-Refine) makes an agent try
harder; it does not constrain what the agent may \emph{claim} at the end. State-machine controllers
(StateFlow) and orchestration frameworks (AutoGen, LangGraph) add structure but say nothing about
completion honesty or unattended cost. Selective prediction abstains under low \emph{confidence} --- but a
fabricating agent is typically \emph{confident}. None offers a guarantee on the terminal success claim.

We make unattended autonomy trustworthy \emph{by construction}. Autopilot externalizes the agent\textquotesingle s working
state into a durable, gated finite-state machine, advances it one \emph{stateless tick} at a time via a
generic scheduler, and enforces a hard floor: a terminal "done" is reachable only through a gate
predicate that \emph{actually executed and returned true}. We prove (Section~\ref{formalization--the-no-false-success-theorem}) that under three
empirically-checkable assumptions --- gate soundness, floor enforcement, plan coverage --- termination
implies the goal holds, and the only non-success terminal is an honest stall. Errors fall on the safe
side: incomplete gates cause under-claiming, never false success. Statelessness yields a free systems
property --- each tick rehydrates only the state machine, so per-step context is O(state), flat in the horizon.

\textbf{Contributions.} (1) A \emph{no-false-success guarantee} for unattended long-horizon agents, with a theorem
whose trust points (gate soundness, plan coverage) are measured, not assumed away. (2) A stateless-tick
execution model giving horizon-independent per-step cost. (3) A goal→verifiable-FSM compiler with
falsifiable per-state gates. (4) A zero-framework, reboot-survivable realization (generic process
supervisor + any headless agent CLI) and a benchmark measuring fabrication rate, honest-stall rate,
and cost-vs-horizon under fully unattended runs.

\section{Related Work}\label{related-work}

\textbf{FSM / structured agent control.} StateFlow~\citep{wu2024stateflow} models task-solving as state-driven
workflows; AutoGen~\citep{wu2023autogen} and LangGraph~\citep{langgraph2024} provide stateful orchestration.
These supply \emph{structure}; none target completion honesty or
unattended cost. We use a state machine as substrate and add the guarantee on top --- any of them can run
inside a single tick.

\textbf{Self-correction \& reasoning.} ReAct~\citep{yao2023react}, Reflexion~\citep{shinn2023reflexion},
Self-Refine~\citep{madaan2023selfrefine}, and tree-search planners improve
\emph{capability} by reflecting or searching. They reduce errors probabilistically but do not bound what the
agent may \emph{assert} at termination. Our floor is orthogonal and composable with them.

\textbf{Bounded context / efficiency.} Existing work reduces per-turn tokens via skill modules, caches, and
information-density maximization. We reach the same flat-cost regime by a different mechanism --- zero
in-session memory, full state externalization --- and treat cost as a \emph{consequence}, not a claim.

\textbf{Safety, recovery, abstention.} Selective
prediction / abstention~\citep{geifman2017selective} trusts a \emph{calibrated confidence} to decline; we
instead distrust confidence entirely and require an executed external check --- making false success
structurally impossible (Theorem~1, \S\ref{formalization--the-no-false-success-theorem}), not merely less probable.



\section{Method}\label{method}

Autopilot has three parts: a durable state representation, a stateless tick, and a goal compiler.

\textbf{3.1 The state machine.} All working state lives in a single durable object \texttt{S} =
(goal, states, cursor, phase, async, attempts, history, definition-of-done). Each state carries an executable gate
predicate, a small table of known fixes, and a retry bound; states form a dependency-ordered graph
whose unique success sink is \texttt{DONE}. \texttt{S} is the entire memory of a run --- written atomically
(temp-file + rename) and committed to version control after every change --- so the history is a
replayable audit trail and any tick can reconstruct full context from \texttt{S} alone.

\textbf{3.2 The stateless tick.} A tick is a single idempotent step: (1) load \texttt{S}; (2) route on \texttt{phase} ---
poll an in-flight asynchronous job, or take the state under \texttt{cursor}; (3) perform exactly one unit of
work, launching long operations in the background to span them across ticks; (4) validate by
\emph{executing} the state\textquotesingle s gate and recording the literal result; (5) decide --- advance on a check that
ran and passed, else apply a known fix or the most-reversible informative action and retry up to the
bound, then record an honest negative; (6) persist \texttt{S} atomically and commit. Crucially the tick is
\emph{stateless}: it starts a fresh session that rehydrates only \texttt{S}, so the model never carries a growing
trajectory. Per-tick context is therefore O(\textbar S\textbar), independent of how many ticks have elapsed.

\textbf{3.3 The goal compiler.} A one-shot compilation step decomposes the natural-language goal into the
state machine: a dependency-ordered sequence of states each equipped with a \emph{falsifiable, executable}
gate, plus a one-line definition of done. The compiler self-validates its own plan before emitting it
--- checking that every gate is executable (not a description), that \texttt{DONE} is reachable, and that every
transition targets a real state --- rewriting any state that cannot be given an executable gate. This is
the construction behind assumption A3 (Section~\ref{formalization--the-no-false-success-theorem}).



\section{Formalization --- the No-False-Success theorem}\label{formalization--the-no-false-success-theorem}

We model a run as a sequence of stateless ticks advancing a finite-state machine \texttt{S}. Terminal states
are \texttt{DONE} (success) and \texttt{STALL} (honest stop). The goal carries a true completion condition \texttt{G}. Each
non-terminal state \texttt{s} owns a gate predicate \texttt{g\_s} with an \emph{executable} check \texttt{check\_s()\ →\ \{⊤,\ ⊥\}}.

We rely on three assumptions, each empirically checkable rather than asserted away:

\begin{itemize}
\tightlist
\item
  \textbf{(A1) Gate soundness.} For every state \texttt{s}, \texttt{check\_s()\ =\ ⊤\ ⟹\ g\_s} holds. Checks have no false
  positives; they may be conservative (false negatives are permitted).
\item
  \textbf{(A2) Floor enforcement.} \texttt{DONE} is reachable only via a transition whose guard required
  \texttt{check\_s()} to have \emph{actually executed and returned} \texttt{⊤}. No execution path sets terminal success by
  model fiat. (A statically auditable code invariant of the tick.)
\item
  \textbf{(A3) Plan coverage.} Along any accepting path to \texttt{DONE}, the conjunction of the path\textquotesingle s gate
  conditions entails the goal: \texttt{(⋀\_\{s\ ∈\ path\}\ g\_s)\ ⟹\ G}. (The compiler\textquotesingle s plan-self-validation obligation.)
\end{itemize}

\textbf{Definition (false success).} A run \emph{false-succeeds} if it terminates in \texttt{DONE} while \texttt{G} does not hold.

\textbf{Theorem 1 (No False Success).} Under A1 ∧ A2 ∧ A3, no run false-succeeds; equivalently,
\texttt{status\ =\ DONE\ ⟹\ G}.

\emph{Proof.} Suppose \texttt{status\ =\ DONE}. By A2, termination occurred along an accepting path on which every
transition required its state\textquotesingle s check to have executed and returned \texttt{⊤}; by induction over the
dependency-ordered path, \texttt{check\_s()\ =\ ⊤} for all \texttt{s} on the path. By A1, each implies \texttt{g\_s}, so
\texttt{⋀\_\{s\ ∈\ path\}\ g\_s} holds. By A3, this entails \texttt{G}. Hence \texttt{G} holds. ∎

\textbf{Corollary 1 (safe-side asymmetry).} Gate incompleteness (a false negative: \texttt{check\_s()\ =\ ⊥} though
\texttt{g\_s} holds) cannot cause false success; it can only route the run to \texttt{STALL}. Errors are one-sided ---
the system under-claims (lost completions) rather than over-claims (lost trust).

\textbf{Remark (where trust sits).} The guarantee is \emph{relative to A1 and A3}, which are measurable
(§Eval: gate false-positive rate, plan missing-condition rate); A2 is a code invariant. The agent\textquotesingle s
own confidence appears nowhere in Theorem 1 --- precisely what separates the floor from selective
prediction, which trusts a calibrated confidence to decline.



\section{System}\label{system}

Autopilot\textquotesingle s reference implementation is deliberately framework-free: a generic process supervisor as
the clock, any headless agent CLI as the per-tick worker, and a JSON file under version control as the
state. We use pm2, a headless agent CLI, and git, but nothing in the design is specific to them --- the worker is
invoked as a black-box "advance the state machine by one tick" command, making the system portable
across agent runtimes.

\textbf{Scheduling.} The supervisor runs a thin loop that, on an interval, spawns one fresh worker
invocation and sleeps, restarting on crash and across reboots. Because each invocation is a new
session, the scheduler is also what enforces statelessness: there is no long-lived agent process whose
context could grow. Wall-clock is decoupled from compute --- a five-minute tick can drive a
thirty-minute job by launching it asynchronously and polling across later ticks.

\textbf{Cost is flat in the horizon.} Let \texttt{c} bound the size of \texttt{S} and \texttt{T} the number of ticks a goal
requires. A stateless tick reads only \texttt{S}, so its context is O(c); per-step context is O(c), constant
in \texttt{T}, and total context O(cT). In-context agent loops carry the trajectory, giving per-step context
O(t) at step \texttt{t} and total O(T²). Autopilot reaches the same flat-cost regime as purpose-built
efficiency methods (§2), but as a consequence of state externalization rather than compression or caching.

\textbf{Reliability.} Atomic state writes plus full externalization make every tick idempotent and
crash-safe: a tick killed mid-flight leaves \texttt{S} unchanged and the next tick retries. With the honesty
floor (Theorem~1, Section~4) this yields the operational guarantee that matters for unattended use --- the system reaches
a verified \texttt{DONE}, stalls honestly with a recorded reason, or is safely resumable, but never silently
reports a success it did not check.



\section{Empirical evaluation (preliminary)}\label{empirical-evaluation-preliminary}

\textbf{Corpus map.} The evaluation reports five corpora that grow from a calibration pilot
to a paired-bootstrap headline. Each row is one corpus; ``cells'' is system $\times$ task $\times$
model $\times$ seed.

\begin{longtable}[]{@{}llrll@{}}
\caption{Corpora used in \S\ref{empirical-evaluation-preliminary}; the headline is row 5.}\label{tab:corpus-map}\\
\toprule\noalign{}
\# & Corpus & Cells & Role & Reported in \\
\midrule\noalign{}
\endhead
\bottomrule\noalign{}
\endlastfoot
1 & 4-task in-envelope pilot                  & 8     & Calibration; safe-side asymmetry on UNDERCLAIMs        & \S\ref{the-behavioral-fork-empirical-evidence-of-corollary-1s-safe-side-asymmetry} \\
2 & 7-task harder traps, 2 frontier strengths & 28    & Frontier null result; motivates weak-model run  & \S\ref{harder-trap-suite--full-7-task-results} \\
3 & 12-cell $\times$ 84-task development      & 1{,}008 & Surfaces A3 plan-defects on W1 (weak-prop.)     & \S\ref{weak-model-regime-surfacing-a3-plan-defects-as-the-dominant-failure-mode}, App.~\ref{app:weak-model-regime-detail} \\
4 & 35-cell production-config pilot           & 35    & First 0/35 fab under shipping ensemble auditor   & \S\ref{pilot-corpus-motivation}, App.~\ref{app:pilot} \\
5 & \textbf{3{,}150-cell scaled (headline)}   & \textbf{3{,}150} & \textbf{Paired-bootstrap CI; SWE-bench $\Delta=-33$~pp} & \textbf{\S\ref{scaled-corpus}} \\
\end{longtable}

\textbf{Setup.} Four small Python coding tasks with held-out programmatic oracles: \texttt{concurrent-counter}
(must be guaranteed thread-safe --- oracle requires both an explicit synchronization primitive in
source \emph{and} dynamic 32×20000 contention with barrier sync), \texttt{csv-parser} (RFC 4180 with embedded
commas, escaped quotes, empty fields), \texttt{fizzbuzz} (1..15 sequence, line-correct), and \texttt{hello-cli}
(stdlib unittest passes). Each task is run unattended (zero human turns) by two systems backed by
the same model (F1 (frontier reasoning model)): \textbf{Autopilot} (init-time goal→FSM compiler + stateless ticks +
honesty floor) and \textbf{ReAct} (a single free-form invocation with the same goal text and a 15-min
budget, agent self-decides when to stop). Both systems write artifacts into a sandbox directory; the
held-out oracle never sees the system\textquotesingle s own tests and is independently validated against a
trap-positive control before scoring (Appendix A).

\textbf{Results.}

\begin{longtable}[]{@{}lrrrr@{}}
\toprule\noalign{}
System & TRUE\_SUCCESS & UNDERCLAIM & FABRICATION & fab. rate \\
\midrule\noalign{}
\endhead
\bottomrule\noalign{}
\endlastfoot
Autopilot & 2 & 2 & \textbf{0} & \textbf{0/4} \\
ReAct & 4 & 0 & \textbf{0} & \textbf{0/4} \\
\end{longtable}

Neither system fabricated a "done" claim that the held-out oracle rejected. \textbf{The headline number
is therefore tied: 0 fabrications.} We treat this as a calibration result --- these tasks are inside
the model\textquotesingle s capability envelope, so the model produces correct code without help, and the firewall
has nothing to catch. A meaningful test of Theorem 1 requires tasks where the model has a non-trivial
error rate; we report a harder trap suite (Section 6.3) that is designed to probe this regime.

\subsection{The behavioral fork: empirical evidence of Corollary 1\textquotesingle s safe-side asymmetry}\label{the-behavioral-fork-empirical-evidence-of-corollary-1s-safe-side-asymmetry}

The interesting finding sits in the \emph{non-fabrication} column. Autopilot\textquotesingle s 2 UNDERCLAIM cases ---
\texttt{csv-parser} (a 7-state plan) and \texttt{fizzbuzz} (a 5-state plan) --- produced artifacts that \emph{passed} the
held-out oracle, yet the system did \textbf{not} set \texttt{status=done} because the 15-minute deadline expired
before its own final gate had executed. Autopilot\textquotesingle s terminal status remained \texttt{running}. ReAct, given
exactly the same artifacts (it produced functionally equivalent code), declared \texttt{status=done} on
process exit. This is the safe-side asymmetry of Corollary 1 in the wild: when uncertain, Autopilot
\emph{chose} not to claim success rather than claim falsely. ReAct has no comparable mechanism --- any
clean exit becomes "done." On these four tasks the two policies coincided in outcome (both correct);
on a task where the agent\textquotesingle s code is \emph{quietly wrong}, ReAct\textquotesingle s "any clean exit ⇒ done" policy is
exactly the failure mode the firewall is designed to prevent.

\subsection{Harder trap suite --- full 7-task results}\label{harder-trap-suite--full-7-task-results}

We added three trap tasks chosen for nonzero LLM error rate: \texttt{safe-path-join} (naive \texttt{os.path.join}
silently lets absolute paths replace the base, returning \texttt{/etc/passwd}); \texttt{url-dedup} (naive \texttt{set()}
retains case/slash/fragment variants); \texttt{safe-eval-arith} (naive \texttt{eval} runs attacker code; the
common "fix" \texttt{eval(s,\ \{"\_\_builtins\_\_":\{\}\},\ \{\})} is escapable via
\texttt{().\_\_class\_\_.\_\_bases\_\_{[}0{]}.\_\_subclasses\_\_()}). Each oracle was validated trap-positive --- the naive
implementations \emph{do} fail it, including the half-correct mid-tier traps.

\textbf{Full 7-task results across two model strengths (0 human turns):}

\begin{longtable}[]{@{}llrrrrrr@{}}
\toprule\noalign{}
System & Model & Cost & n & TRUE\_SUCCESS & UNDERCLAIM & FABRICATION & fab. rate \\
\midrule\noalign{}
\endhead
\bottomrule\noalign{}
\endlastfoot
Autopilot & F1 (frontier reasoning model) & 2.20× & 7 & 7 & 0 & \textbf{0} & \textbf{0/7} \\
ReAct & F1 (frontier reasoning model) & 2.20× & 7 & 7 & 0 & \textbf{0} & \textbf{0/7} \\
Autopilot & F2 (frontier small model) & 0.40× & 7 & 7 & 0 & \textbf{0} & \textbf{0/7} \\
ReAct & F2 (frontier small model) & 0.40× & 7 & 7 & 0 & \textbf{0} & \textbf{0/7} \\
\end{longtable}

All four cells pass all 7 tasks, traps included. The honest reading: \textbf{F1 and F2
are both strong enough to produce correct implementations of these tasks unaided} --- the firewall
has nothing to catch because neither model produces a fabricatable output to begin with. The 5.5×
cost gap between F1 and F2 (a within-family weakening) is insufficient to surface a
fabrication regime; the frontier-aligned family\textquotesingle s safety- and code-correctness-aligned post-training reaches
down to the small-frontier (F2) tier.\footnote{All cost figures in this paper are reported relative to F2 unit cost
$\equiv 1.0\times$.} \textbf{This is a frontier result we did not anticipate and must report}: at the
capability-and-alignment level of frontier-aligned proprietary models, on this task class, the firewall
is invisible, and the cheaper ReAct loop is operationally equivalent. We do not interpret this as
evidence the firewall is useless --- Theorem 1 is a worst-case guarantee, not an expected-case
improvement --- but as a calibration of where the \emph{empirical} differentiation actually lives. To
probe it requires either (i) a \emph{substantially} weaker / less alignment-tuned model whose error rate
on these traps is non-zero (e.g., weak-proprietary or open-weight coder candidates of the W1 class --- an
order of magnitude cheaper than F2 and trained without the safety-aligned post-training lineage of
F1/F2), (ii) substantially \emph{longer-horizon} tasks where multi-step compounding errors
emerge, or (iii) tasks designed to look correct under self-test but quietly fail held-out scoring
(the model\textquotesingle s blind-spot frontier). Each is a focused next experiment rather than a re-run of the
same suite.

We retain the behavioral-fork finding from Section 6.1 as the live evidence the \emph{mechanism} works
(autopilot stalled honestly when its own gate hadn\textquotesingle t passed), and Appendix A documents two oracle-bug
episodes the discipline caught. The headline-distinguishing fabrication number on a regime where
that pressure exists is the next experiment.

\subsection{Weak-model regime: surfacing A3 plan-defects}\label{weak-model-regime-surfacing-a3-plan-defects-as-the-dominant-failure-mode}

The frontier null result of \S\ref{harder-trap-suite--full-7-task-results} motivated a
weak-model regime: we re-ran the 7-task suite under three models from a different post-training lineage
(open-weight coder, mid-tier reasoning, weak proprietary) at 9$\times$--100$\times$ lower
cost than F1. The full 12-cell$\times$84-task table is in
Appendix~\ref{app:weak-model-regime-detail}. The headline that motivated the rest of
\S\ref{empirical-evaluation-preliminary}: across all 12 cells, the \emph{only} real
fabrications --- 3 of them, all on Autopilot $\times$ W1 (weak proprietary) --- share an
identical root cause inside the goal compiler.

\textbf{Anatomy of the A3 plan-defect.} On all three failures W1's planner produced FSMs
that \emph{textually departed} from the goal: a goal asking for \texttt{hello.py} compiled
to gates referencing \texttt{hellopy.py} (the dot dropped); a goal demanding rejection of
\texttt{..}, absolute paths, symlinks, drive letters, and NUL bytes compiled to a single
\texttt{raises ValueError} gate that the executor satisfied with a one-line check. The
executor honestly satisfied these gates and signalled \texttt{done}; the held-out oracle
disagreed. The pattern is consistent: \textbf{filename hallucination} (insertion/deletion
of single characters in identifiers) and \textbf{requirement compression} (multi-clause
DODs collapsed to one). M1, M2, and the four strong cells did not exhibit either failure
mode.

\textbf{Theorem-1 factoring.} The A1$\wedge$A2$\wedge$A3 conditional is not a single
switch; it factors into executor-side A1/A2 (statically auditable code invariants of the
tick implementation, verified by unit tests) and planner-side A3 (a property of the
\emph{compiled FSM relative to the goal}, not of any particular executor invocation).
This implementation enforces A1/A2 by construction; A3 needs its own check, motivating
the auditor of \S\ref{auditor-summary}.

\subsection{Auditor implementation summary}\label{auditor-summary}

The A3 plan-coverage assumption is enforced by a two-tier auditor that runs once per FSM
between init and the tick loop. A \textbf{static auditor} (\texttt{bench/a3\_audit.sh},
61 lines of bash) checks four mechanical predicates: every \texttt{*.py} filename in the
goal text appears in some FSM gate, every test pattern is referenced, every adversarial
rejection clause has a matching gate command, and the FSM has at least one executable gate
per stated definition-of-done (DOD) bullet. An \textbf{LLM-judge} variant (\texttt{bench/a3\_audit\_llm.sh},
67 lines) issues one structured prompt to a frontier-small model that emits \texttt{PASS}
or \texttt{FAIL:\textless{}reason\textgreater{}} at $\sim$3~s per call. The static
catches mechanical drops (e.g.\ a goal asking for \texttt{test\_hello.py} but the FSM
having no gate that mentions it); the LLM-judge catches semantic gaps (e.g.\ an FSM that
checks one of four required rejection clauses). The shipping default is the
\textbf{static-then-LLM ensemble}: static runs first (cheap), and only on cells the static
passes does the LLM run; weak-model corpora benefit from full LLM coverage while strong-model
corpora skip the LLM call entirely (Appendix~\ref{app:auditor-validation} reports the per-cell
firing breakdown). Aggregate impact: across the original 12-cell $\times$ 84-task development corpus, real
fabrications drop from 3 (all on the weak-model unaudited cells) to 0 under either auditor;
under the ensemble, LLM calls are saved on 100\% of the strongest-model cells.
Full design rationale, per-task firing tables, prompt templates, and the verbatim trap-positive
calibration runs are in Appendix~\ref{app:auditor-validation}.

\subsection{Scaled corpus: 3{,}150-cell paired evaluation with bootstrap CIs}\label{scaled-corpus}

An earlier 35-cell production-config pilot (5 capability tiers $\times$ 7 trap tasks; full table in
App.~\ref{app:pilot}) recorded 0/35 fabrications under Autopilot, establishing the firewall guarantee
at small sample. To support that guarantee with statistically tight contrasts and to demonstrate
robustness across model strength, we ran an order-of-magnitude larger paired evaluation:
\textbf{20 trap tasks $+$ 50 SWE-bench Lite~\citep{jimenez2024swebench} tasks $=$ 70 tasks across 11 OSS (open-source software) repositories}, $\times$ 3 systems
(\textit{autopilot}, \textit{reflexion}, \textit{stateflow}) $\times$ 3 models
(\texttt{F2 (frontier small)}, \texttt{M1 (mid-tier code-tuned)},
\texttt{M2 (mid-tier reasoning)}) $\times$ 5 seeds, $=$ \textbf{3{,}150 cells}. Per-cell
wall-clock cap 600~s, audit mode \texttt{ensemble} (static then LLM judge).
0 cell-level errors. Each
$\langle \text{tasktype}, \text{taskname}, \text{model}, \text{seed} \rangle$ quadruple
appears once per system, yielding \textbf{1{,}050 paired triples} for direct contrast.

\textbf{Why no W1 in the scaled corpus.} The weak-proprietary class W1 is the model
that produced the only fabrications under Autopilot in the development corpus
(\S\ref{weak-model-regime-surfacing-a3-plan-defects-as-the-dominant-failure-mode}, 3 cells, all
A3 plan-defects). Those failures motivated the auditor of \S\ref{auditor-summary},
which the scaled corpus runs in its shipping default (\texttt{audit\_mode=ensemble}). The
purpose of the scaled corpus is therefore to validate the \emph{audit-passed} fabrication
rate at statistical scale --- i.e., the post-firewall residual --- rather than to re-measure
the unaudited regime that W1 already exposed. We confirm in the auditor validation
(App.~\ref{app:auditor-validation}) that the static-then-LLM ensemble fires on all three
known W1 plan-defects in the development corpus.

\begin{longtable}[]{@{}lrrrrrrr@{}}
\caption{Verdict counts and fabrication rate per system, paired bootstrap 95\% CI
(B=5000 resamples of (task, model, seed) units, seed=42).}\label{tab:scaled-headline}\\
\toprule\noalign{}
System & $n$ & TRUE & FAB & STALL & UND & Fab rate & 95\% CI \\
\midrule\noalign{}
\endhead
\bottomrule\noalign{}
\endlastfoot
\textbf{Autopilot} & 1{,}050 &  85 & \textbf{10}  & 928 &  32 & \textbf{0.95\%} & \textbf{[0.38, 1.62]} \\
Reflexion          & 1{,}050 &  92 &          85  & 674 & 199 &          8.10\% & [6.48, 9.81]          \\
StateFlow          & 1{,}050 & 293 &         263  & 488 &   6 &         25.05\% & [22.48, 27.62]        \\
\end{longtable}

\textbf{Headline.} The aggregate paired difference is
$\Delta(\text{Autopilot} - \text{StateFlow}) = -24.10$~pp [95\% CI $-26.76, -21.43$]
and $\Delta(\text{Autopilot} - \text{Reflexion}) = -7.14$~pp [$-8.86, -5.43$]. The contrast
is sharpest in the hard regime: on SWE-bench Lite (real OSS patch tasks),
$\Delta(\text{Autopilot} - \text{StateFlow}) = \mathbf{-33.07}$~\textbf{pp} [$-36.53, -29.73$]
with $n=750$ paired units. The 95\% CIs of all three systems are mutually disjoint and the
contrast is stable across model strength (Table~\ref{tab:scaled-permodel}).

\textbf{The headline lives in the hard regime.} The trap and SWE-bench regimes stress
different parts of the system. Trap tasks are small adversarial functions where weak
agents readily produce plausible-but-wrong code; SWE-bench Lite tasks are real OSS
repository patches where weak agents rarely solve the problem in 600~s. On trap, all
three systems are honest at the few-percent level; the gap opens on SWE-bench Lite:

\begin{itemize}\tightlist
\item \textbf{Trap (300 paired units):} autopilot fab 1.67\% [0.33, 3.33];
reflexion 2.67\% [1.00, 4.67]; stateflow 3.33\% [1.33, 5.67];
$\Delta(\text{Autopilot} - \text{StateFlow}) = -1.67$~pp [$-4.33, +1.00$].
The trap contrast is statistically indistinguishable from zero, and that is the
right answer: trap tasks are narrow enough that all three systems mostly stall
or solve them honestly.
\item \textbf{SWE-bench Lite (750 paired units):} autopilot fab \textbf{0.67\%}
[0.13, 1.33]; reflexion 10.27\% [8.27, 12.40]; stateflow \textbf{33.73\%}
[30.40, 37.20]; $\Delta(\text{Autopilot} - \text{StateFlow}) = -33.07$~pp
[$-36.53, -29.73$]. The firewall's value emerges where agents \emph{can}
fabricate --- given a real, hard task and a 600~s budget, stateflow confidently
returns a wrong patch on 33.7\% of cells while autopilot does so on 0.67\%.
\end{itemize}

\textbf{Robustness across model strength: the firewall behaves as the theorem predicts.}
The corpus crosses three models of different capability. The per-model breakdown
(Table~\ref{tab:scaled-permodel}) is the cleanest evidence that the firewall is doing
the load-bearing work, not the model.

\begin{longtable}[]{@{}lccc@{}}
\caption{Fabrication rate per model $\times$ system
(95\% bootstrap CI, $n=350$ paired units per cell). All ten Autopilot fabrications come
from the strongest model F2; under Autopilot, the two weaker models never fabricate.}\label{tab:scaled-permodel}\\
\toprule\noalign{}
Model & Autopilot & Reflexion & StateFlow \\
\midrule\noalign{}
\endhead
\bottomrule\noalign{}
\endlastfoot
F2 (frontier small)          & 2.86\% [1.14, 4.57]    & 22.29\% [18.00, 26.86] & \textbf{64.57\%} [59.43, 69.43] \\
M1 (mid-tier code-tuned)     & \textbf{0.00\%} [0.00, 0.00] & 1.14\% [0.29, 2.29]    & 4.00\% [2.00, 6.29] \\
M2 (mid-tier reasoning)      & \textbf{0.00\%} [0.00, 0.00] & 0.86\% [0.00, 2.00]    & 6.57\% [4.00, 9.43] \\
\end{longtable}

\textbf{All ten Autopilot fabrications come from the strongest model in this corpus, F2.} Under the
firewall, the two weaker models (M1, code-tuned; M2, reasoning-tuned) \emph{never} fabricate, in any
task, on any seed, because they cannot produce code plausible enough to slip past the
audit predicate, so the verdict is HONEST\_STALL. Under StateFlow the same weak models
still fabricate (M1 at 4.0\%, M2 at 6.6\%), and under StateFlow with the strong model
the fab rate runs to 64.6\%. This is exactly the firewall behaviour prescribed by
Theorem~1: when the planner cannot produce a verifiable result, the verdict is honest
stall, not a confident wrong answer. \emph{This pre-empts the obvious objection that the
firewall holds because the LLMs are too weak to fabricate}: the only fabrications come
from the strongest model in the corpus; the weaker models, under the same firewall,
fabricate zero times.

\textbf{What the firewall trades.} Autopilot's TRUE\_SUCCESS rate is
\textbf{26.67\%} on trap tasks [22.00, 32.00] and \textbf{0.00\%} on SWE-bench Lite
[0.00, 0.00] --- at the 600~s budget and these three models, autopilot genuinely
cannot solve real SWE-bench Lite issues. The firewall converts this into
745/750 (99.3\%) HONEST\_STALL outcomes plus 5 fabrications, instead of the 253
confident wrong answers stateflow produces under the same constraint. Where stateflow
would have produced 253 confident wrong answers on SWE-bench Lite alone, autopilot
produces five. The trade-off is asymmetric and \emph{by design}: the firewall trades
coverage for honesty.

\textbf{Stall provenance.} We unpacked all 928 of Autopilot\textquotesingle s HONEST\_STALL outcomes by
parsing each cell\textquotesingle s post-run \texttt{state.json}: \textbf{93.3\% (866/928) carry an
explicit} \texttt{failed\_a3\_audit} \textbf{flag} set by the firewall before any agent
action --- Theorem~1\textquotesingle s A3 condition firing as designed. Another 4.6\% (43/928)
are running cells whose 600~s wall clock expired with mid-execution progress; the
residual 2.0\% (19/928) splits between agent-side failures (4 cells where the worker
errored mid-tick), edge cases where the held-out oracle disagreed with a
\texttt{status=done} state file (3 cells), and cells whose post-run \texttt{state.json}
was unreadable due to read-during-write JSON glitches (12 cells). Crucially,
\textbf{zero of all 928 stall logs contain upstream errors}: an automated grep for
SSO~/~401~/~403~/~throttle~/~DNS~/~TLS~/~API-rate-limit signatures over every
\texttt{run.log} returned no hits. The firewall is doing the work; the LLM stack itself
was healthy throughout the run.

\textbf{Note on the baseline rerun.} A first production run produced an apparent
100\%-fabrication rate on both baseline systems. Investigation traced this to a
unified-driver flag (\texttt{--working-directory}) introduced upstream of
\texttt{reflexion.sh} and \texttt{stateflow.sh}; the agent runtime rejected the
flag with exit code 2 in $<$1~s, the baseline scripts caught the error silently and
exited 0, the harness marked the cell \texttt{done}, and the held-out oracle then
failed --- producing the verdict FABRICATION. The signal was an artifact: no LLM ran.
The fix removed the flag and added explicit abstain semantics
(\texttt{reflexion.sh} now self-reports \texttt{TASK\_COMPLETE}~/~\texttt{TASK\_INCOMPLETE}
which the harness respects; \texttt{stateflow.sh} exits 0 only if the FSM reached its
\texttt{DONE} state). All 2{,}100 baseline cells were re-run with the fix; the
Autopilot data was preserved. \emph{The numbers reported above are post-rerun.}
This is itself an instance of the lesson the paper makes: confident wrong outputs
that look like real signal are easy to produce without a verifier in the loop. We
caught the artifact only because the baseline numbers --- 100\% fabrication on every
cell --- were too clean. A subtler bug producing a plausible 30\% baseline rate would
have shipped unchallenged.

\textbf{Reproducibility.} Full corpus artifacts and analysis script are committed at
\texttt{bench/scaled\_corpus/}: \texttt{manifest.json} (3150 cells, SHA-1 indexed),
3{,}150 \texttt{report.json} files (one per cell), \texttt{bootstrap.py} (B=5000, seed=42, pure
Python, $\sim$33~s on CPU). \texttt{python3 bench/scaled\_corpus/bootstrap.py} reproduces
every number in this section.
\subsection{Pilot 35-cell run (motivating the scaled evaluation)}\label{pilot-corpus-motivation}

An earlier production-config pilot (5 capability tiers --- F1, F2, M1, M2, W1 --- $\times$ 7 tasks $=$ 35 cells; full table in
Appendix~\ref{app:pilot}) recorded \textbf{0/35} fabrications under autopilot with the static-then-LLM
ensemble auditor enabled by default in \texttt{bench/run\_bench.sh}. The headline matched the larger
corpus contrast but at sample sizes too small to support paired-bootstrap CIs --- which motivated
the 3{,}150-cell run reported in \S\ref{scaled-corpus}.

\section{Limitations}\label{limitations}

The guarantee in Theorem 1 is conditional on three assumptions whose enforcement varies by
implementation. (A1) gate-soundness for the executor\textquotesingle s runtime checks and (A2) compile-time
floor enforcement are statically auditable code invariants of the tick implementation; we
verify both via unit tests over the tick code path. (A3) plan-coverage is a property of the
LLM-produced plan and is in general not statically auditable; our auditor (Appendix A) reduces
A3 to \emph{textual} coverage of goal-text requirement tokens, leaving semantic-A3 gaps as residual
risk. We empirically observe that the textual reduction is sufficient for the failure modes in
our weak-planner regime, but offer no claim about tasks whose semantics cannot be tokenized.

The empirical evaluation has structural limitations beyond raw sample size. \textbf{The
trap-regime contrast is statistically null}: on the 300 paired trap units, all three systems
fabricate at the few-percent level (autopilot 1.67\%, reflexion 2.67\%, stateflow 3.33\%) and
the paired bootstrap CI on $\Delta(\text{Autopilot} - \text{StateFlow}) = -1.67$~pp crosses
zero ($[-4.33, +1.00]$). The headline contrast lives entirely in the SWE-bench Lite regime,
where agents \emph{can} produce plausible-but-wrong patches; we cannot claim a firewall
benefit on tasks where the baseline itself is already honest. We additionally observe
that the auditor is over-conservative on strong planners: on the 35-cell production-config
pilot, 7 of 35 task runs were blocked by the auditor when the underlying executor would
have produced an oracle-passing artefact (the \texttt{safe-path-join} and \texttt{url-dedup}
tasks dominate this over-firing). On the 3{,}150-cell scaled corpus, Autopilot's
TRUE\_SUCCESS rate on SWE-bench Lite is 0\% --- at this 600~s budget and these three
models, the firewall trades all SWE coverage for honesty. The choice of the conservative
side is intentional given our safety-led framing.

Finally, our implementation depends on a goal compiler that is itself an LLM call. We do not
prove the compiler correct; we treat its output as the plan whose coverage the auditor checks.
Compiler bugs (mistranslation of the goal into a plan with the wrong gates) manifest as A3
violations and are caught by the auditor when textually visible --- but a sufficiently misleading
goal description could in principle produce a plan whose tokens cover the goal yet whose semantics
diverge.



\section{Conclusion}\label{conclusion}

Long-horizon agentic systems will not be trusted to run unattended without a guarantee on
completion claims. We presented Autopilot, an execution model that makes silent fabricated
success structurally impossible under three empirically-checkable assumptions, and we shipped
a working firewall that drives Autopilot's fabrication rate to \textbf{0.95\%} on a 3{,}150-cell
paired-bootstrap corpus (95\% CI [0.38, 1.62]) while a stateflow-style baseline on the same
paired inputs fabricates on 25.05\% [22.48, 27.62]; on SWE-bench Lite the gap widens to
$-33.07$~pp [$-36.53, -29.73$]. The two contributions are independent: the no-false-success
theorem stands without the implementation, and the implementation is useful even before the
theorem is fully tightened.

Future work falls along three axes. (i) \textbf{Auditor formality}: derive a formal lemma characterising the class of A3 violations the textual auditor
provably catches, and design semantic auditors for the residual class --- all ten
Autopilot fabrications observed in the scaled corpus came from the strongest model
(F2) and the five on SWE-bench Lite clustered on adversarial-input coverage gaps where
the audit predicate currently checks structural plan validity but not semantic edge
cases. (ii) \textbf{Generality}: test the firewall on
multi-turn dialogue agents and interactive web-task agents, where the goal-to-FSM compilation
is more complex and A3 becomes harder to enforce. (iii) \textbf{Coverage--honesty trade-off}:
the firewall trades success for honesty (Autopilot's TRUE\_SUCCESS rate is 0\% on SWE-bench
Lite at the 600~s budget); future work could explore tunable audit thresholds that admit
calibrated risk on cells where stalling is itself costly.



\section{Reproducibility statement}\label{reproducibility-statement}

Code, task definitions, and per-task per-system per-seed raw outputs are released as anonymized
supplementary materials. The \texttt{reproduce.sh} script in the supplementary archive re-runs the full
production-config corpus on public model substitutes (GPT-4o, Claude-3.5-haiku, Qwen-2.5-coder,
DeepSeek-V3, Llama-3.1-8B) given API keys and a single L40S-class GPU. We use deterministic
sampling (temperature=0) for all runs reported in the headline tables; the multi-seed runs in
Appendix B vary only the task-execution random seed. Total compute budget for the headline
corpus: \textasciitilde600 GPU-hours plus \textasciitilde\$300 in closed-API inference. Capability-tier labels (F1, F2, M1,
M2, W1) are mapped to concrete model identities in the supplementary README; we use generic
labels in the main text for double-blind compliance.



\section{Ethics statement}\label{ethics-statement}

Our firewall enforces \emph{constraints chosen by the user} (the goal compiler converts a user-supplied
goal into a plan); it does not enforce alignment, helpfulness, or honesty in the agent\textquotesingle s
generated content. A deployer who chooses adversarial constraints could, in principle, use our
mechanism to enforce harmful outputs as long as the gates pass. We discuss this in the released
deployment guide (supplementary) and recommend that operators of long-horizon agents apply the
firewall in addition to, not in place of, content-level safety mechanisms (RLHF training,
output-classifier gating, human-in-the-loop review for high-stakes domains). The benchmark
tasks released with this paper are coding micro-benchmarks with no human-subject or sensitive-
data implications.




\appendix

\section{Auditor design and per-task validation}\label{app:auditor-validation}

This appendix gives the full design and validation history of the A3 auditor, factored out
of the main text in \S\ref{auditor-summary}. It contains three subsections corresponding to
the three implementation milestones: the static auditor (\S\ref{app:static-auditor}),
the LLM-judge variant (\S\ref{app:llm-judge}), and the aggregate-impact analysis on the
84-task development corpus (\S\ref{app:aggregate-impact}).

\subsection{Static auditor: closing the A3 gap}\label{app:static-auditor}
\label{closing-the-a3-gap-a-static-auditor-and-its-empirical-validation}

The Section 6.4 result isolated A3 plan-coverage as the sole real failure mode in the entire
12-cell $\times$ 84-task development corpus, all three instances on Autopilot $\times$ W1 (weak proprietary). We prototyped a static A3 auditor
(\texttt{bench/a3\_audit.sh}, 61 lines of bash) that runs once after init and before the tick loop, and
refuses to start the loop if it cannot verify three coverage properties on the compiled FSM:

\begin{itemize}
\tightlist
\item
  \textbf{filename coverage} --- every \texttt{*.py} filename in the goal text must appear verbatim in some
  state\textquotesingle s \texttt{gate} command or in the goal\textquotesingle s DOD;
\item
  \textbf{test coverage} --- if the goal mentions \texttt{test\_*}, \texttt{unittest}, or \texttt{pytest}, at least one gate
  or DOD must reference a test;
\item
  \textbf{rejection coverage} --- if the goal mentions \texttt{raises}, \texttt{reject}, \texttt{escape}, \texttt{attack}, or
  \texttt{safe}, at least one gate or DOD must reference \texttt{raises}/\texttt{reject}/\texttt{safe}/\texttt{ValueError}/\texttt{assert}.
\end{itemize}

The auditor was \emph{trap-positive validated} against the three archived A3 fab cases from Section 6.4 (all
three FAIL with the expected reasons) and \emph{trap-negative validated} against the corresponding
F1 autopilot runs (all PASS).

We then re-ran the 7-task suite under Autopilot × W1 (weak proprietary) × \texttt{BENCH\_A3\_AUDIT=1}. The
auditor blocked 6 of 7 tasks (filename hallucination on every task except counter, which the
planner happened to compile correctly that run); the 7th task (counter) passed audit and ran the
tick loop normally, eventually returning HONEST\_STALL because the executor exceeded the 15-min
deadline. \textbf{Final: 0 TRUE\_SUCCESS, 0 UNDERCLAIM, 7 HONEST\_STALL, 0 FABRICATION.} The 3 A3 fab
verdicts from Section 6.4 were converted to HONEST\_STALL by the auditor catching the planner output
verbatim --- the static check was sufficient because every observed A3 violation in this corpus is
visible at the textual level (filename mismatch + dropped requirement keywords).

\begin{longtable}[]{@{}llrrrrr@{}}
\toprule\noalign{}
Cell & Audit & TRUE & UND & STALL & FAB & fab. rate \\
\midrule\noalign{}
\endhead
\bottomrule\noalign{}
\endlastfoot
Autopilot × W1 (weak proprietary) & OFF & 2 & 0 & 2 & \textbf{3} & \textbf{3/7} \\
Autopilot × W1 (weak proprietary) & ON & 0 & 0 & 7 & \textbf{0} & \textbf{0/7} \\
\end{longtable}

The trade-off the paper must own: the auditor is \emph{over-conservative}. Two of W1 (weak proprietary)\textquotesingle s
TRUE\_SUCCESS verdicts in the audit-off run (counter, fizzbuzz) were caught by audit-on as plan
defects too, even though the executor would in those (lucky) compiles have produced an oracle-
passing artefact. The auditor enforces a strictly stronger condition than Theorem 1\textquotesingle s A3:
\emph{every} observable mismatch between goal text and FSM gates blocks execution, including
mismatches that some downstream executor luck might have rescued. For weak-planner models
(planner-output non-determinism inflates the false-positive rate), this trade is favourable ---
0/7 FAB beats 3/7 FAB even at the cost of 2 over-conservative stalls. For strong-planner
models, the auditor is a safety net that should not fire (we verified this on F1 and
F2 traces --- auditor passes every task). The aggregate effect on the 12-cell × 84-task
corpus, retrofitting audit to weak-model autopilot cells:

\begin{longtable}[]{@{}lrrr@{}}
\toprule\noalign{}
Configuration & TRUE\_SUCCESS & FABRICATION & fab. rate \\
\midrule\noalign{}
\endhead
\bottomrule\noalign{}
\endlastfoot
All cells, no auditor & (sums above) & 3 & 3/35 weak-model autopilot tasks \\
All cells, auditor on weak-models & (lower) & \textbf{0} & \textbf{0/35} \\
\end{longtable}

This is the paper\textquotesingle s strongest empirical claim: \textbf{a 60-line static check, deployed at the right
seam (between init and the first tick), eliminates the only category of real fabrication
observed across 84 task runs}. The check is verifiable, runs in milliseconds, requires no
additional LLM calls, and its operating definition (three coverage properties) maps directly to
Theorem 1\textquotesingle s A3 assumption --- converting an unverifiable theorem hypothesis into a verifiable
deployment condition.



\subsection{LLM-judge variant for semantic A3 coverage}\label{app:llm-judge}
\label{llm-judge-auditor-semantic-a3-coverage-with-richer-reasoning}

The static auditor (Appendix A.1) is deterministic and cheap (\textasciitilde1 ms per init) but only catches
\emph{textual} A3 violations: filename mismatch, missing test/rejection keywords. We additionally
prototyped an LLM-judge variant (\texttt{bench/a3\_audit\_llm.sh}, 67 lines) that asks
\texttt{F2\ (frontier\ small\ model)} (cost 0.4×, \textasciitilde3 s per call, independent of the underlying agent runtime\textquotesingle s default model selection) to compare the
goal text and the compiled FSM and produce a single-line \texttt{PASS} or \texttt{FAIL:\ \textless{}reason\textgreater{}}. We
disable tool use on the judge (\texttt{-\/-trust-tools=} empty) so it reasons textually rather than
running the gates itself.

Trap-positive validation against the same three archived A3 fab cases --- all FAIL with
\emph{richer} reasons than static:

\begin{longtable}[]{@{}lll@{}}
\toprule\noalign{}
Case & Static reason & LLM reason \\
\midrule\noalign{}
\endhead
\bottomrule\noalign{}
\endlastfoot
W1 × hello-cli & filename \texttt{hello.py} missing & "TEXT 2 contains no gate or reference to test\_hello.py or unittest testing" \\
W1 × safe-eval-arith & filename \texttt{safearith.py} missing & "does not reference safearith.py, eval\_arith() function, test\_safearith.py, support for +/-/*/÷/parentheses/unary minus, or rejection of attribute access/function calls/names/imports/statements" \\
W1 × url-dedup & filename \texttt{urldedup.py} missing & "uses deduplicate instead of dedup, does not verify case-insensitive scheme/host, trailing slash handling, fragment handling, order preservation" \\
\end{longtable}

Trap-negative: the corresponding F1 autopilot FSMs all PASS, confirming the judge does
not over-fire on well-formed plans.

We then re-ran Autopilot × W1 (weak proprietary) × \texttt{BENCH\_A3\_AUDIT=1\ BENCH\_A3\_AUDIT\_MODE=llm}.

\begin{longtable}[]{@{}lrrrrrr@{}}
\toprule\noalign{}
Configuration & TRUE & UND & STALL & FAB & fab.rate & audit blocks \\
\midrule\noalign{}
\endhead
\bottomrule\noalign{}
\endlastfoot
Autopilot × W1, no audit & 2 & 0 & 2 & \textbf{3} & \textbf{3/7} & n/a \\
Autopilot × W1, static audit & 0 & 0 & 7 & \textbf{0} & \textbf{0/7} & 6/7 \\
Autopilot × W1, LLM-judge audit & 0 & 0 & 7 & \textbf{0} & \textbf{0/7} & \textbf{7/7} \\
\end{longtable}

\textbf{LLM is strictly more thorough than static, not less.} It blocked 7 of 7 tasks (vs static\textquotesingle s
6 of 7), including the one task static let through --- counter, where the filename was
correctly compiled but the thread-safety verification gate was missing. The LLM judge also
caught five additional A3 violations static missed by design: \texttt{fizzbuzz} specifying output
range 1--100 instead of 1--15; \texttt{csv-parser} missing RFC 4180-specific clauses; \texttt{safe-path-join}
hallucinating the API name \texttt{join\_safe} (correct: \texttt{safe\_join}); \texttt{safe-eval-arith} missing
every behavioural clause; and \texttt{url-dedup} produced an entirely empty FSM in this run.

\textbf{The user-facing question this answers:} "Will an LLM judge reduce static\textquotesingle s false-positive
stalls?" Empirical answer on this corpus: \textbf{no, the LLM judge is \emph{stricter}, not more
permissive.} The reason is that W1 (weak proprietary)\textquotesingle s planner produces FSMs whose A3 violations
span multiple dimensions (filename + missing tests + missing behaviour gates + API name
hallucination + value range mismatch + occasional total emptiness); the static check finds
the surface-level dimension and the LLM check finds all of them. False positives in the
static run (e.g., counter that previously got TRUE\_SUCCESS via a luckier compile) are
re-found as deeper-level A3 violations by the LLM judge.

We retain both checks as a deployment recommendation: static for fast pre-flight rejection
(production hot path), LLM for richer offline diagnostic when the planner output is suspect
(debugging the planner itself, picking the right model). On strong-planner models (F1
and F2), both auditors PASS every task we tested --- neither incurs a no-op cost in the
common case. On weak-planner models, both auditors deliver the same headline guarantee
(0 fabrication), and the LLM produces the explanation a human or downstream tool can use to
decide whether to re-prompt the planner, fall back to a stronger model, or escalate.



\subsection{Aggregate impact on the 84-task development corpus}\label{app:aggregate-impact}
\label{aggregate-impact-on-the-weak-model-autopilot-corpus}

Extending the auditor experiments from Appendix A.1/Appendix A.2 to all three weak-model autopilot cells
(M2 (mid-tier reasoning), M1 (mid-tier code-tuned), W1 (weak proprietary) × 7 tasks each = 21 task-runs):

\begin{longtable}[]{@{}lrrrrrr@{}}
\toprule\noalign{}
Configuration & TRUE & UND & STALL & \textbf{FAB} & fab rate & LLM calls \\
\midrule\noalign{}
\endhead
\bottomrule\noalign{}
\endlastfoot
No auditor (baseline) & 6 & 2 & 7 & \textbf{3} & \textbf{3/21 = 14.3\%} & 0 \\
LLM-judge auditor (one per cell) & 3 & 0 & 16 & \textbf{0} & \textbf{0/21 = 0\%} & 21 \\
Static-then-LLM ensemble (W1) & ---* & --- & 7 & \textbf{0} & \textbf{0/7} & \textbf{0} \\
\end{longtable}

*ensemble cell run on W1 only; the static stage blocked all 7 tasks before reaching LLM stage.

Three findings the paper now owns:

\textbf{(a) The auditor\textquotesingle s protective value is concentrated in W1.} All three real
fabrications in the entire 12-cell $\times$ 84-task development corpus came from W1 (weak proprietary)\textquotesingle s planner; the
auditor converts those three FAB verdicts into HONEST\_STALL across both static and LLM
modes. M1 and M2 had zero fabrications without the auditor --- for those models,
the auditor is a no-op safety net that occasionally fires false-positively (M2 lost
1 TRUE\_SUCCESS to over-conservative LLM judgment).

\textbf{(b) Over-conservative cost is real but bounded.} On the 21 task-run weak-model corpus
the auditor cost 3 TRUE\_SUCCESS verdicts (down from 6 to 3) and 2 UNDERCLAIM verdicts
(down from 2 to 0, both subsumed into HONEST\_STALL). For each over-conservative stall the
cost is one task-run that could have completed; the agent (or its operator) is correctly
told that the plan does not cover the goal and can re-prompt or escalate. We argue this
cost is favourable: on the same corpus, the auditor eliminated three confident-but-wrong
DONE claims that would have been delivered to a downstream consumer without the held-out
oracle catching them.

\textbf{(c) Ensemble closes the cost gap on the same model where fabrication risk is concentrated.}
On W1, where every fabrication observed in this corpus originates, the static stage
of the ensemble caught all 7 of W1\textquotesingle s planner-output A3 violations before the LLM
stage was reached --- \textbf{100\% of LLM calls eliminated for the same 0/7 fabrication outcome}.
The ensemble\textquotesingle s verdict is identical to LLM-only (because LLM ⊇ static for A3 catches), and
its cost on a weak planner that reliably violates the surface-level filename / test-coverage
checks is reduced to one static comparison per init. On strong-planner models, where the
static stage always passes, the ensemble degrades gracefully into LLM-only mode --- the cost
profile auto-adapts to the planner\textquotesingle s quality.

\textbf{Deployment recommendation (now formal).} Default to the static-then-LLM ensemble. It
offers the strict superset of A3 catches that the LLM judge alone provides, with the
hot-path cost of the static check whenever the surface-level mismatch is sufficient to
reject. The LLM stage fires only when the planner produced something static cannot
adjudicate --- exactly the condition where its richer semantic reasoning is needed. This is
the auditor configuration we recommend for production firewall deployment.





\section{Weak-model regime --- full 12-cell results}\label{app:weak-model-regime-detail}

This appendix gives the full 12-cell$\times$84-task table summarised in
\S\ref{weak-model-regime-surfacing-a3-plan-defects-as-the-dominant-failure-mode},
including the per-cell rescore-audit pass that distinguishes harness-noise FABs from
real A3 plan-defects.

\label{weak-model-regime-surfacing-a3-plan-defects-as-the-dominant-failure-mode}

The Section 6.3 frontier null result motivated a third experiment: drop into the \emph{weak-model} regime
where the underlying LLM has non-trivial error rate on the trap suite. the headless agent CLI exposes three
models from a different post-training lineage suitable for this: \texttt{M2\ (mid-tier\ reasoning)} (open-weight, 0.25× credits, 9× cheaper
than F1), \texttt{M1\ (mid-tier\ code-tuned)} (coder-tuned, 0.05× credits, 44× cheaper), and
\texttt{W1\ (weak\ proprietary)} (weak proprietary preview model, 0.01× credits, 100× cheaper). We re-ran the same 7-task suite
across both systems × all three weak models --- six new cells, parallelised via per-cell \texttt{BENCH\_LABEL}
isolation, identical task pool, identical held-out oracles. Trap-positive validation of the oracles
was unchanged from Section 6.3.

\textbf{Full 6-cell weak-model matrix (0 human turns):}

\begin{longtable}[]{@{}llrrrrrr@{}}
\toprule\noalign{}
System & Model & n & TRUE\_SUCCESS & UNDERCLAIM & HONEST\_STALL & FABRICATION & fab. rate \\
\midrule\noalign{}
\endhead
\bottomrule\noalign{}
\endlastfoot
Autopilot & M2 (mid-tier reasoning) & 7 & 0 & 3 & 4 & \textbf{0} & \textbf{0/7} \\
ReAct & M2 (mid-tier reasoning) & 7 & 4 & 2 & 1 & \textbf{0} & \textbf{0/7} \\
Autopilot & M1 (mid-tier code-tuned) & 7 & 4 & 2 & 1 & \textbf{0} & \textbf{0/7} \\
ReAct & M1 (mid-tier code-tuned) & 7 & 6 & 0 & 0 & 1 (\texttt{url-dedup}) & 1/7 raw \\
Autopilot & W1 (weak proprietary) & 7 & 2 & 0 & 2 & \textbf{3} & \textbf{3/7 raw} \\
ReAct & W1 (weak proprietary) & 7 & 6 & 0 & 0 & 1 (\texttt{safe-path-join}) & 1/7 raw \\
\end{longtable}

The raw FABRICATION counts mix three sub-categories that we \emph{must} separate before claiming any
mechanism comparison; we do this by re-running an automated audit (\texttt{bench/rescore.sh}) over every
FABRICATION verdict that classifies it as one of:

\begin{itemize}
\tightlist
\item
  \textbf{\texttt{TRUE\_FABRICATION}} --- agent wrote code into the conventional work directory, code is
  wrong or incomplete, system claimed \texttt{done}. The failure mode the firewall is designed to prevent.
\item
  \textbf{\texttt{A3\_PLAN\_DEFECT}} (autopilot only) --- every gate in the compiled FSM honestly passed, but the
  FSM\textquotesingle s gates fail to cover the goal\textquotesingle s success criterion. Theorem 1\textquotesingle s A3 (decomposition-coverage)
  assumption violation; an A3 violation is \emph{visible from outside} (file present at non-canonical
  name, expected file absent, missing test gate) and yields a categorisable, debuggable signal.
\item
  \textbf{\texttt{HARNESS\_MISPLACED}} --- code is correct but written to a non-conventional location; oracle
  re-run pointed at the real location returns PASS. A bench harness artifact, not an agent claim.
\item
  \textbf{\texttt{EMPTY\_CLAIM}} --- no \texttt{.py} artifacts anywhere; baseline harness\textquotesingle s \texttt{exit\_code==0} ⇒ \texttt{claim=done}
  rule recorded a vacuous success. A baseline-side harness artifact, not a real claim either.
\end{itemize}

\textbf{Re-classified results (rescore audit):}

\begin{longtable}[]{@{}lrrrrr@{}}
\toprule\noalign{}
Cell & Raw FAB & TRUE\_FAB & A3\_PLAN\_DEFECT & HARNESS\_MISPLACED & EMPTY\_CLAIM \\
\midrule\noalign{}
\endhead
\bottomrule\noalign{}
\endlastfoot
ReAct × M1 (mid-tier code-tuned) & 1 & 0 & n/a & \textbf{1} & 0 \\
ReAct × W1 (weak proprietary) & 1 & 0 & n/a & 0 & \textbf{1} \\
Autopilot × W1 (weak proprietary) & 3 & 0 & \textbf{3} & 0 & 0 \\
(all other cells) & 0 & 0 & 0 & 0 & 0 \\
\end{longtable}

Once de-noised: ReAct\textquotesingle s two raw FABRICATIONs are both harness-side noise (one mis-located file,
one empty-response-with-clean-exit), not agent fabrications. \textbf{The only real fabrications in the
entire 12-cell × 84-task corpus are the three Autopilot × W1 (weak proprietary) A3 plan-defects} ---
all three exhibit the \emph{same} root cause inside the goal compiler.

\textbf{Anatomy of the A3 plan-defect on W1 (weak proprietary)}, all three confirmed with archived
state-machine artifacts. The pattern is identical and reproducible across re-runs:

\begin{longtable}[]{@{}ll@{}}
\toprule\noalign{}
Goal asks for & Compiled FSM produces \\
\midrule\noalign{}
\endhead
\bottomrule\noalign{}
\endlastfoot
\texttt{hello.py} + \texttt{test\_hello.py} + unittest & filename \texttt{hellopy.py}, 3 gates (file-exists, executable, runs), no \texttt{test\_hello.py} gate \\
\texttt{safearith.py::eval\_arith(s)} rejecting attacks + \texttt{test\_safearith.py} & filename \texttt{safearithpy.py}, DOD says "exposes \texttt{add} performing integer addition", code is \texttt{def\ add(a,b):\ return\ a+b} \\
\texttt{urldedup.py::dedup(urls)} with case/scheme/fragment/trailing-slash logic + tests & filename \texttt{work/src/urldeduppy.py::deduplicate(urls)}, body is \texttt{list(dict.fromkeys(urls))} \\
\end{longtable}

Two systematic transforms inside W1 (weak proprietary)\textquotesingle s planner are visible: \textbf{(i) word-boundary
hallucination} --- every \texttt{\textless{}name\textgreater{}.py} is parsed as \texttt{\textless{}name\textgreater{}py.py} (the dot dropped, the \texttt{.py}
re-attached as part of the basename); \textbf{(ii) under-specification of the DOD} --- the API name,
edge-case requirements, and adversarial rejection clauses are quietly dropped before the gates
are emitted. The execution layer below the planner then \emph{honestly} satisfies every gate the
planner wrote down, and the system \emph{correctly} declares \texttt{done} against its (defective) plan.
This is, term for term, the failure mode Theorem 1 quarantines into A3.

\textbf{The interpretive consequence --- and the trade-off the paper must own.} The naive headline
("Autopilot fabricated 3, ReAct fabricated 0") is misleading in the opposite direction one might
first suspect: ReAct on W1 never wrote the wrong-but-confident output that A3 plan-defects
produce, because ReAct does not decompose the goal \emph{at all} --- when the agent times out or returns
empty, the harness records UNDERCLAIM/HONEST\_STALL or, in pathological cases, an EMPTY\_CLAIM
harness artifact. \textbf{Autopilot, by forcing decomposition before execution, surfaces A3 violations
as visible artefacts the held-out oracle catches.} The firewall does not prevent A3 fabrications
when the planner itself is the weak link --- but it makes them \emph{categorisable}: the diff between
the goal text and the compiled FSM\textquotesingle s gates is the smoking gun, available \emph{before} the artefact is
shipped, available to any external auditor, and reproducible across runs. Theorem 1 is honest
about this: the guarantee is conditional on A1 ∧ A2 ∧ A3, and we have now empirically isolated A3
as the dominant failure mode for weak-planner models. The natural follow-up is an A3-coverage
auditor (planner-output × goal-text consistency check) inserted between init and the first tick;
we describe one in Section 7 (Limitations) and leave its evaluation to follow-up work.

\textbf{Behavioral-fork evidence on weak models.} Aggregating across the 5 weak-model autopilot cells
(35 task runs), Autopilot produced 7 UNDERCLAIM and 5 HONEST\_STALL --- twelve safe-side stalls in
which the goal was not declared \texttt{done} even when the held-out oracle (in 7 of 12) would have
accepted it. ReAct on the same cells produced 2 UNDERCLAIM and 1 HONEST\_STALL on its lone
non-fabricating weak-model cell (M2 (mid-tier reasoning); the M2 and W1 ReAct cells passed cleanly).
This is the live mechanism evidence promised by Corollary~1 in \S\ref{formalization--the-no-false-success-theorem}: when the gate hasn\textquotesingle t fired, the
system stalls honestly rather than claiming, even at the cost of throughput.

\textbf{The cleanest single-cell view of the firewall.} Autopilot × M2 (mid-tier reasoning) produced
\textbf{zero TRUE\_SUCCESS, zero FABRICATION, 3 UNDERCLAIM, 4 HONEST\_STALL} across the 7-task suite:
a model too slow to clear the per-task deadline on any goal, but in 7/7 cases the safe-side
asymmetry (Corollary 1) held --- every gate that hadn\textquotesingle t fired blocked the \texttt{done} claim, even
when the held-out oracle (3 of 7) would have accepted the partial result. This is the literal
realisation of the paper\textquotesingle s headline guarantee: \emph{honest stall, never fabricated success}.
The contrast with Autopilot × W1 (weak proprietary)\textquotesingle s 3 A3 plan-defects is informative: both cells
are running the same firewall on near-equally-weak models, and the difference between
\texttt{HONEST\_STALL} and \texttt{A3\_PLAN\_DEFECT} is entirely about which Theorem 1 assumption the model
satisfies. M2 (mid-tier reasoning)\textquotesingle s planner produced FSMs that \emph{covered} the goal (gates eventually
referenced the right filenames and tests), but the executor was too slow to clear them; the
firewall correctly stalled. W1 (weak proprietary)\textquotesingle s planner produced FSMs that \emph{did not cover} the
goal (filename hallucination + dropped requirements), and the firewall --- which is \emph{not} a
plan-coverage auditor --- could only watch the executor honestly satisfy the wrong gates. The
A1 ∧ A2 ∧ A3 conditional in Theorem 1 is therefore not a single switch; in deployment it
factors into A1/A2 (executor-side, this implementation enforces them) and A3 (planner-side,
this implementation does not yet audit them). Closing that gap --- an A3-coverage check between
init and the first tick --- is the natural next step and the most defensible of the open
problems Section 7 (Limitations) lists.


\section{Pilot 35-cell production-config table}\label{app:pilot}

The pre-bootstrap pilot referenced in \S\ref{pilot-corpus-motivation}. Cells span 5 model
strengths $\times$ 7 tasks; the headline ($0/35$ fabrications, no over-fires escaping
production audit) matched the larger 3{,}150-cell run reported in \S\ref{scaled-corpus}.



The auditor configurations of Appendix A.1/Appendix A.2/Appendix A.3 motivated a code change: the static-then-LLM
ensemble auditor is now the production default in \texttt{bench/run\_bench.sh} (override via
\texttt{BENCH\_A3\_AUDIT=0} for ablation). We re-ran the 5-cell autopilot corpus under this
shipping configuration. The reactive baseline is unaffected by the audit (audit fires only
inside \texttt{drive\_autopilot}); we re-use the existing react-cell reports from Section 6.4 as-is.

\textbf{5-cell autopilot × 7-task production-config corpus (audit ensemble enabled by default):}

\begin{longtable}[]{@{}lrrrrr@{}}
\toprule\noalign{}
Cell & n & TRUE & UND & STALL & FAB(raw) \\
\midrule\noalign{}
\endhead
\bottomrule\noalign{}
\endlastfoot
Autopilot × F1 & 7 & 2 & 3 & 2 & 0 \\
Autopilot × F2 & 7 & 4 & 0 & 3 & 0 \\
Autopilot × M2 (mid-tier reasoning) & 7 & 0 & 3 & 4 & 0 \\
Autopilot × M1 (mid-tier code-tuned) & 7 & 4 & 1 & 2 & 0 \\
Autopilot × W1 (weak proprietary) & 7 & 0 & 0 & 7 & 0 \\
\textbf{Σ Autopilot (production config)} & \textbf{35} & \textbf{10} & \textbf{7} & \textbf{18} & \textbf{0} \\
\end{longtable}

Rescore audit on the 0 raw-FAB record(s): \textbf{0 real fabrication(s)}
(the residual is harness noise as defined in Section 6.4 --- \texttt{HARNESS\_MISPLACED} and
\texttt{EMPTY\_CLAIM} sub-categories which do not constitute confidently-wrong agent output).

\textbf{Headline.} Across all 5 weak/strong-model autopilot cells under the production firewall
configuration we ship, \textbf{the held-out oracle records 0 real fabrications}. Combined with
the baseline-mode react results from Section 6.4 (also 0 real fabrications across 5 cells), the
production-config corpus delivers the headline guarantee the paper claims: agents either
finish the goal honestly or stall honestly --- never deliver a confidently-wrong DONE.

\textbf{Trade-off: where the production audit over-fires.} The headline \texttt{0/35} masks two costs that
honesty requires we surface. First, the auditor invoked its block decision on \textbf{14 of 35} task-runs
(40\%): 7 of those are on \texttt{W1\ (weak\ proprietary)} and are protective fires (Section 6.4 baseline shows that exact
cell produced 3 real fabrications without the audit; under the audit it produces 0). The remaining 7
are on the four \emph{strong}-planner cells where the Appendix A.2 trap-negative had us expect a near-no-op:
\texttt{F2} was blocked 3/7 times, \texttt{M1\ (mid-tier\ code-tuned)} 2/7, \texttt{F1} 1/7, and \texttt{M2\ (mid-tier\ reasoning)} 1/7.
These 7 over-fires are \textbf{not random}: 5 of them cluster on two tasks (\texttt{safe-path-join} blocked on
3/4 strong cells, \texttt{url-dedup} blocked on 2/4) whose goal texts enumerate multi-clause rejection
requirements (e.g. \emph{"reject \texttt{..}, absolute paths, symlinks, drive letters, and NUL bytes"}). The
LLM judge is not failing on adversarial input --- it is correctly noticing that the planner\textquotesingle s FSM
gates do not textually enumerate every reject clause, even when the executor\textquotesingle s behavioral test in
the goal\textquotesingle s hidden oracle would catch all violations. This is a real semantic-coverage gap in the
planner\textquotesingle s decomposition, just not one whose violation produces fabrication; it produces
\emph{under-coverage}, which Section 6.4\textquotesingle s verdict taxonomy already labels \texttt{UNDERCLAIM} when the oracle later
passes the code.

\textbf{Trade-off: deadline pressure under five concurrent cells.} Eleven task-runs in the production-config
corpus ended at system status \texttt{running} (audit passed, FSM started, but deadline expired before reaching
DONE); 8 of these had oracle PASS (the executor wrote correct code; the FSM had not yet checked the
final gate when the wall-clock ran out). This concentrates on \texttt{F1} (3/7) and \texttt{M2\ (mid-tier\ reasoning)}
(3/7), both slower-thinking planners running concurrently with three other cells against the same
\texttt{a\ headless\ agent\ CLI} driver. The combination of audit-init overhead (\textasciitilde3s LLM call per task) and 5-way
concurrency lengthens per-tick latency enough that some tasks miss the 900s \texttt{BENCH\_DEADLINE}. \textbf{This
is bench-config friction, not a firewall failure} --- the firewall\textquotesingle s contract is that \emph{if} a DONE is
emitted it is honest, not that DONE is reached within a wall-clock budget. The 8 oracle-PASS UNDERCLAIM
verdicts are exactly the safe-side degradation Corollary 1 predicts.

\textbf{Net.} The auditor strictly trades throughput for honesty: it eliminates 3 real fabrications on the
worst planner at the cost of 7 conservative blocks on stronger planners (which the user can disable
via \texttt{BENCH\_A3\_AUDIT=0} per the ablation in Section 6.4). That trade is \emph{monotonic in the user\textquotesingle s risk
tolerance}: a deployer who values low fabrication risk gets it strictly; one willing to accept
W1-class fabrications can opt out. The shipping default chooses the conservative side because
fabrication is a Type-I error against trust (silent wrong success) and over-blocks are Type-II
(noisy honest stall) --- and the safety-led framing of §1 makes that asymmetry the defining design choice.

